% !TeX root = ../../Book5.tex
% 纲要标题：### 4.2 Mackey 分解
% 纲要位置：工作记录/计划纲要.md，第 4144 行。
% 纲要细目（写作范围）：
% * 双陪集公式
% * 交织算子空间的计算
% * 不可约性判据


\section{Mackey 分解}
\label{b5:sec:0402}

双陪集分块中的 \(x\otimes W\) 受交集 \(K\cap xHx^{-1}\) 作用。这个作用不是原封不动地把 \(W\) 限制到交集：必须先通过 \(x\) 把交集搬回 \(H\)。将这一点写清以后，Mackey 公式便可由一个显式张量映射得到。

\subsection{双陪集公式}
\label{b5:subsec:0402:01}

\begin{definition}\label{b5:def:conjugate-representation}
	设 \(G\) 为有限群，\(H\leq G\)，\(k\) 为域，\(\rho:H\to\operatorname{GL}_k(W)\) 为有限维表示，\(x\in G\)。在同一向量空间 \(W\) 上定义 \(xHx^{-1}\) 的作用
	\[
		\rho^x(a)=\rho(x^{-1}ax)\qquad(a\in xHx^{-1}).
	\]
	所得表示记为 \({}^xW\)，称为 \(W\) 的\term[gonge-biaoshi]{共轭表示}[conjugate representation]。
\end{definition}
\symidx{\({}^xW\)：共轭表示}

这个方向约定与左陪集相配合：在诱导表示中， \(a(x\otimes w)=x\otimes\rho(x^{-1}ax)w\)。此外，先取 \({}^yW\) 再取其 \(x\)-共轭，所得表示为 \({}^{xy}W\)。

\begin{theorem}[Mackey 分解]\label{b5:thm:mackey-decomposition}
	设 \(G\) 为有限群，\(H,K\leq G\)，\(k\) 为任意域，\(W\) 为有限维左 \(k[H]\)-模。取 \(K\backslash G/H\) 的一组代表元 \(X\)，并令 \(L_x=K\cap xHx^{-1}\)。则有左 \(k[K]\)-模同构
	\begin{equation}\label{b5:eq:mackey-decomposition}
		\operatorname{Res}_K^G\operatorname{Ind}_H^G W \cong \bigoplus_{x\in X} \operatorname{Ind}_{L_x}^K \operatorname{Res}_{L_x}^{xHx^{-1}}({}^xW).
	\end{equation}
\end{theorem}
这一分解称为 \term[mackey-fenjie]{Mackey 分解}[Mackey decomposition]，公式也称为 Mackey 双陪集公式。
\begin{proof}
	根据\cref{b5:prop:double-coset-blocks}，只须识别每一块。记 \(W\) 的 \(H\)-作用为 \(\rho\)。固定 \(x\)，把 \({}^xW\) 限制到 \(L_x\)，定义
	\[
		\begin{aligned}
			\Phi_x:k[K]\otimes_{k[L_x]}{}^xW &\longrightarrow k[KxH]\otimes_{k[H]}W,\\
			a\otimes w&\longmapsto ax\otimes w.
		\end{aligned}
	\]
	对 \(\ell\in L_x\)，有
	\[
		a\ell x\otimes w =ax(x^{-1}\ell x)\otimes w =ax\otimes\rho(x^{-1}\ell x)w.
	\]
	这正是源端平衡关系的像，故映射良定义；左 \(K\)-线性直接来自左乘。它的逆映射可以在群元素基上写成
	\[
		\Psi_x(k_0xh\otimes w)=k_0\otimes hw \qquad(k_0\in K,\ h\in H).
	\]
	这里须检查同一群元素的不同写法。如果 \(k_0xh=k_1xh_1\)，则
	\[
		\ell=k_1^{-1}k_0\in K\cap xHx^{-1}=L_x, \qquad (x^{-1}\ell x)h=h_1.
	\]
	源端的平衡关系因而给出
	\[
		k_0\otimes hw =k_1\ell\otimes hw =k_1\otimes\rho(x^{-1}\ell x)hw =k_1\otimes h_1w.
	\]
	所以 \(\Psi_x\) 与所选写法无关。对于 \(h_0\in H\)，它还把 \(k_0xhh_0\otimes w\) 与 \(k_0xh\otimes h_0w\) 都送到 \(k_0\otimes hh_0w\)，故经线性延拓后确实定义了张量积上的映射。直接代入可得 \(\Psi_x\Phi_x=1\)、\(\Phi_x\Psi_x=1\)。交集 \(L_x\) 正是同一元素写成 \(kxh\) 时产生的歧义所在，源端对 \(L_x\) 取张量积恰好消除了这种歧义。对所有 \(x\) 取直和即得结论。
\end{proof}

选取别的代表元会改变上式各项的具体写法，却不改变\cref{b5:prop:double-coset-blocks}中由双陪集本身确定的子空间。具体地，若 \(x'=k_0xh_0\)，其中 \(k_0\in K\)、\(h_0\in H\)，则 \(L_{x'}=k_0L_xk_0^{-1}\)，且对 \(\ell\in L_x\) 有
\[
	\begin{aligned}
		\rho^{x'}(k_0\ell k_0^{-1}) &=\rho(h_0^{-1}x^{-1}\ell xh_0),\\
		\rho(h_0)\rho^{x'}(k_0\ell k_0^{-1}) &=\rho^x(\ell)\rho(h_0).
	\end{aligned}
\]
后一等式使映射
\[
	k[K]\otimes_{k[L_{x'}]}{}^{x'}W \longrightarrow k[K]\otimes_{k[L_x]}{}^xW, \qquad a\otimes w\longmapsto ak_0\otimes\rho(h_0)w
\]
与平衡关系相容；其逆为 \(a\otimes w\mapsto ak_0^{-1}\otimes\rho(h_0^{-1})w\)。再经 \(\Phi_x\)，像为 \(ak_0x\otimes h_0w=ax'\otimes w\)，与 \(\Phi_{x'}\) 相同。因此两个代表元确实识别了同一个双陪集块。

\ProseParagraphBreak

回到 \(S_3\)、\(H=K=\langle(12)\rangle\) 的符号表示 \(W\)。代表元 \(1\) 的交集为 \(H\)，对应的项就是 \(W\)；代表元 \((13)\) 的交集为平凡群，对应的项是 \(\operatorname{Ind}_{\{1\}}^H\mathbb C=\mathbb C[H]\)。于是 Mackey 公式给出 \(W\oplus\mathbb C[H]\cong\mathbf1_H\oplus W^{\oplus2}\)，与上节直接计算的两块作用一致。

\subsection{交织算子空间的计算}
\label{b5:subsec:0402:02}

\begin{theorem}[交织空间公式]\label{b5:thm:mackey-intertwining}
	设 \(G\) 为有限群，\(H,K\leq G\)，\(k\) 为域，\(U,W\) 分别为有限维左 \(k[K]\)-模、左 \(k[H]\)-模。取双陪集集 \(K\backslash G/H\) 的代表元集 \(X\)，对 \(x\in X\) 令 \(L_x=K\cap xHx^{-1}\)。有 \(k\) 向量空间同构
	\begin{equation}\label{b5:eq:mackey-intertwining}
		\begin{aligned}
			&\operatorname{Hom}_{k[G]} \left(\operatorname{Ind}_K^G U,\operatorname{Ind}_H^G W\right)\\
			&\qquad\cong \bigoplus_{x\in X}\operatorname{Hom}_{k[L_x]} \bigl(\operatorname{Res}_{L_x}^K U, \operatorname{Res}_{L_x}^{xHx^{-1}}({}^xW)\bigr).
		\end{aligned}
	\end{equation}
\end{theorem}
\begin{proof}
	\cref{b5:thm:frobenius-reciprocity}把左侧识别为 \(\operatorname{Hom}_{k[K]}(U,\operatorname{Res}_K^G\operatorname{Ind}_H^G W)\)。 将\cref{b5:thm:mackey-decomposition}代入第二个变量，便得到各个
	\[
		\operatorname{Hom}_{k[K]} \bigl(U,\operatorname{Ind}_{L_x}^K \operatorname{Res}_{L_x}^{xHx^{-1}}({}^xW)\bigr)
	\]
	的直和。

	此时诱导出现在 Hom 的第二个变量，不能再次直接使用同一方向的左伴随。\cref{b5:thm:induction-coinduction}说明有限群的诱导与余诱导自然同构；\cref{b5:prop:coinduction-adjunction}给出的右伴随性质把这一项识别为
	\[
		\operatorname{Hom}_{k[L_x]} \bigl(\operatorname{Res}_{L_x}^K U, \operatorname{Res}_{L_x}^{xHx^{-1}}({}^xW)\bigr).
	\]
\end{proof}

在复数域上，取特征标后，公式也计算两个诱导特征标的内积。取 \(K=H\)、\(U=W=\mathbf1_H\)，每一项都是一维空间，故
\[
	\dim_{\mathbb C}\operatorname{End}_{\mathbb C[G]} \bigl(\mathbb C[G/H]\bigr)=|H\backslash G/H|.
\]
这个等式中每个双陪集对应的算子也可以直接写出。对双陪集 \(D=HxH\)，在陪集基 \(e_{gH}\) 上定义
\[
	T_D(e_{gH})=\sum_{yH\subseteq gD}e_{yH}.
\]
\(D\) 对右乘 \(H\) 稳定，所以 \(gD\) 是若干左 \(H\)-陪集之并，求和有意义；\(D\) 对左乘 \(H\) 稳定，所以 \((gh)D=gD\)，公式与 \(gH\) 所选代表元无关。对 \(a\in G\)，左乘把 \(gD\) 中的陪集逐一送到 \(agD\) 中的陪集，故 \(T_D(ae_{gH})=aT_D(e_{gH})\)。不同 \(D\) 的算子作用于 \(e_H\) 后得到支集互不相交的非零向量，因此这些算子线性无关；由上述维数等式，它们组成自同态空间的一组基。它的矩阵正是 \(G/H\times G/H\) 上轨道 \(\{(yH,gH):yH\subseteq gD\}\) 的示性矩阵，与\cref{b5:prop:permutation-character}中的轨道矩阵一致。

把 \(e_{gH}\) 识别为 \(G/H\) 上的点示性函数，就可用函数卷积记录这些算子的乘法。设 \(f,g:G\to\mathbb C\) 都在左右乘 \(H\) 下不变，\(\varphi:G/H\to\mathbb C\)，规定
\begin{equation}\label{b5:eq:finite-hecke-convolution-action}
	\begin{aligned}
		(f*g)(z)&=\frac1{|H|}\sum_{y\in G}f(y)g(y^{-1}z),\\
		(A_f\varphi)(xH)&=\frac1{|H|}\sum_{y\in G}f(x^{-1}y)\varphi(yH).
	\end{aligned}
\end{equation}
双边不变性保证第一式仍为双边不变函数，并使第二式与 \(xH\) 的代表元无关。左平移作用为 \((a\varphi)(xH)=\varphi(a^{-1}xH)\)，在求和中把 \(y\) 改为 \(ay\)，便知它与 \(A_f\) 交换。两个算子的复合满足
\[
	\begin{aligned}
		(A_fA_g\varphi)(xH) &=\frac1{|H|^2}\sum_{y,z\in G}f(x^{-1}y)g(y^{-1}z)\varphi(zH)\\
		&=\frac1{|H|}\sum_{z\in G}(f*g)(x^{-1}z)\varphi(zH),
	\end{aligned}
\]
故 \(A_fA_g=A_{f*g}\)，即卷积的次序与算子复合的次序相同。单位函数是 \(1_H\)，因为 \(A_{1_H}=I\)。按上述陪集基算子的定义，\(A_{1_D}=T_{D^{-1}}\)：其作用于 \(e_{gH}\) 后，在 \(xH\) 处非零的条件是 \(x^{-1}g\in D\)，亦即 \(xH\subseteq gD^{-1}\)。双陪集示性函数给出全部双边不变函数的一组基，而这些算子给出自同态空间的一组基，因此这个卷积代数与 \(\operatorname{End}_G\mathbb C[G/H]\) 相互识别。

\ProseParagraphBreak

在 \(H=\langle(12)\rangle\) 时的 \(S_3/H\) 三点模型中，\(T_H=I\)，而另一个双陪集的算子把每个基向量送到其余两个基向量之和。若 \(J\) 表示所有矩阵元均为 \(1\) 的三阶矩阵，则后一算子就是 \(T=J-I\)，从而自同态空间为 \(\mathbb CI\oplus\mathbb CT\)。两个双陪集均与其逆相同，所以在函数模型中 \(A_{1_H}=I\)、\(A_{1_D}=T\)。由 \(J^2=3J\)，有
\[
	T^2=(J-I)^2=J+I=T+2I, \qquad 1_D*1_D=1_D+2\cdot1_H.
\]
这里 \(|H|^{-1}\) 的归一化使 \(1_H\) 成为单位，也使非平凡双陪集的示性函数恰对应 \(J-I\)，不需要再乘 \(|H|\)。双陪集的计数在这里变成了两类具体的交织算子及其乘法。

\ProseParagraphBreak

例如，\(S_n\) 在 \(n\) 个点上的置换表示对应 \(H=S_{n-1}\)；当 \(n\geq2\) 时只有两个双陪集，所以它的自同态空间维数为 \(2\)。平凡表示出现一次，其余直和项的重数平方和因而为 \(1\)，得到一个不可约的 \((n-1)\) 维分量。

\subsection{不可约性判据}
\label{b5:subsec:0402:03}

\begin{theorem}[Mackey 不可约性判据]\label{b5:thm:mackey-irreducibility}
	设 \(G\) 为有限群，\(H\leq G\)，\(W\) 为非零有限维复 \(H\)-表示。则 \(\operatorname{Ind}_H^G W\) 不可约，当且仅当 \(W\) 不可约，且对每个不等于 \(H\) 的双陪集 \(HxH\)，令 \(L_x=H\cap xHx^{-1}\) 后有
	\[
		\operatorname{Hom}_{\mathbb C[L_x]} \bigl(\operatorname{Res}_{L_x}^H W, \operatorname{Res}_{L_x}^{xHx^{-1}}({}^xW)\bigr)=0.
	\]
\end{theorem}
\begin{proof}
	\cref{b5:thm:maschke}说明有限群的复表示完全可约。若一个非零表示分解为互不同构不可约表示 \(V_i\) 的 \(\bigoplus_iV_i^{\oplus m_i}\)，则\cref{b5:prop:isotypic-decomposition}给出其自同态空间维数为 \(\sum_i m_i^2\)；这个数等于 \(1\)，恰好意味着表示不可约。

	对\cref{b5:thm:mackey-intertwining}取 \(K=H\)、\(U=W\)，代表元 \(x=1\) 的项就是 \(\operatorname{End}_{\mathbb C[H]}W\)。所有直和项的维数非负，且这一项至少为 \(1\)。因此总维数为 \(1\)，当且仅当这一项维数为 \(1\) 而其余项都为零。再次使用完全可约性与上述维数公式，前一条件等价于 \(W\) 不可约。
\end{proof}

代表元 \(1\) 的项记录 \(W\) 原有的自同态，其余双陪集则检查 \(W\) 与共轭表示在交集上的限制之间是否还存在非零交织。只要存在这样一项，诱导表示就有额外的自同态；对有限群的复表示，完全可约性和上述重数公式说明，自同态空间维数超过 \(1\) 便迫使表示可约。例如取 \(G=S_3\)、\(H=\langle(12)\rangle\)，对非平凡双陪集仍有 \(L_{(13)}=\{1\}\)。若 \(W\) 是任意非零有限维复 \(H\)-表示，令 \(d=\dim_{\mathbb C}W\)，则这一项是两个 \(d\) 维空间之间的全部线性映射，维数为 \(d^2\)。因此
\[
	\dim_{\mathbb C}\operatorname{End}_{\mathbb C[S_3]} \bigl(\operatorname{Ind}_H^{S_3}W\bigr) \geq 1+d^2>1,
\]
所以从这个非正规子群诱导的非零表示都可约。对于前面的符号表示，代表元 \(x=(13)\) 对应的交织项维数为 \(1\)，加上单位双陪集的一维项，自同态空间的维数为 \(2\)。不可约重数的平方和因此为 \(2\)，这与\cref{b5:exm:s3-transposition-induction}中 \(\operatorname{Ind}_H^{S_3}W\cong\varepsilon\oplus U\) 中两个互不同构且各出现一次的不可约分量一致。

\ProseParagraphBreak

不可约性判据中的基域条件不能从前面的模同构公式中自动删去。在模特征下，表示可能不是完全可约的；此时自同态空间一维并不足以推出不可约。

\begin{example}\label{b5:exm:dihedral-induction}
	设 \(G=D_{2n}=\langle r,s:r^n=s^2=1,\ srs=r^{-1}\rangle\)，其中 \(n\geq3\)，且 \(H=\langle r\rangle\)。令 \(\zeta\) 为本原 \(n\) 次单位根，取 \(\lambda_j(r)=\zeta^j\)。因为 \(H\) 正规，双陪集代表元为 \(1,s\)，交集均为 \(H\)，并且 \({}^s\lambda_j=\lambda_{-j}\)。所以 \(\operatorname{Ind}_H^G\lambda_j\) 不可约，当且仅当 \(2j\not\equiv0\pmod n\)。

	\ProseParagraphBreak
	在基 \(1\otimes1,s\otimes1\) 下，两个生成元的矩阵为
	\[
		r\longmapsto
		\begin{pmatrix}\zeta^j&0\\0&\zeta^{-j}\end{pmatrix},
		\qquad s\longmapsto
		\begin{pmatrix}0&1\\1&0\end{pmatrix}.
	\]
	当 \(2j\equiv0\) 时，向量 \((1,1)\) 和 \((1,-1)\) 各自生成不变直线，诱导表示分成两个一维表示。这也从矩阵上解释了判据中的例外。
\end{example}
